\documentclass[11pt]{article}
\newcommand{\DocNumber}{VOLUME 08}
\newcommand{\DocTitle}{C++23 Prerequisite-Based Technical Curriculum}
\newcommand{\DocSubtitle}{Twenty-eight-week technical learning program}
\newcommand{\DocSummary}{Sixteen dependent learning modules, measurable outcomes, labs, assessment gates and a secure capstone.}
% Shared, pdfLaTeX-only design system for the C++23 documentation suite.
\usepackage[T1]{fontenc}
\usepackage[utf8]{inputenc}
\usepackage{lmodern}
\usepackage[english]{babel}
\usepackage[letterpaper,margin=0.82in,headheight=16pt,headsep=17pt]{geometry}
\usepackage{microtype}
\usepackage{amsmath,amssymb}
\usepackage{array,booktabs,tabularx,longtable}
\usepackage[table]{xcolor}
\usepackage[most]{tcolorbox}
\usepackage{listings}
\usepackage{fancyhdr}
\usepackage{titlesec}
\usepackage{enumitem}
\usepackage{hyperref}
\usepackage{needspace}
\usepackage{etoolbox}
\definecolor{ink}{HTML}{13253B}
\definecolor{teal}{HTML}{006E7F}
\definecolor{muted}{HTML}{5A6C7C}
\definecolor{fog}{HTML}{EDF4F6}
\definecolor{sand}{HTML}{FAF3E8}
\definecolor{rule}{HTML}{C9D7DF}
\definecolor{codebg}{HTML}{F7F9FB}
\definecolor{codecomment}{HTML}{47715D}
\hypersetup{colorlinks=true,linkcolor=teal,urlcolor=teal,pdfauthor={Independent C++23 engineering documentation},pdftitle={C++23 Technical Suite}}
\setlength{\parindent}{0pt}
\setlength{\parskip}{6pt}
\setcounter{tocdepth}{2}
\renewcommand{\arraystretch}{1.18}
\titleformat{\section}{\Large\bfseries\color{ink}}{\thesection}{0.7em}{}
\titleformat{\subsection}{\large\bfseries\color{teal}}{\thesubsection}{0.55em}{}
\titlespacing*{\section}{0pt}{17pt}{7pt}
\titlespacing*{\subsection}{0pt}{11pt}{4pt}
\pagestyle{fancy}
\fancyhf{}
\fancyhead[L]{\footnotesize\bfseries\color{teal}C++23 / ISO 14882:2024}
\fancyhead[R]{\footnotesize\color{muted}\DocNumber}
\fancyfoot[L]{\scriptsize\color{muted}Independent engineering companion, not an ISO publication}
\fancyfoot[R]{\scriptsize\color{muted}\thepage}
\renewcommand{\headrulewidth}{0.35pt}
\lstdefinestyle{cppstyle}{language=C++,basicstyle=\ttfamily\footnotesize,keywordstyle=\color{teal}\bfseries,commentstyle=\color{codecomment},stringstyle=\color{ink},backgroundcolor=\color{codebg},frame=single,rulecolor=\color{rule},breaklines=true,breakatwhitespace=true,columns=fullflexible,keepspaces=true,showstringspaces=false,tabsize=2,numbers=none,upquote=true,xleftmargin=3pt,xrightmargin=3pt,aboveskip=6pt,belowskip=8pt}
\lstset{style=cppstyle}
\newtcolorbox{standardbox}{enhanced,breakable,colback=fog,colframe=teal!35!white,boxrule=0.55pt,arc=1.4mm,left=3mm,right=3mm,top=2mm,bottom=2mm,before skip=7pt,after skip=8pt}
\newtcolorbox{practicebox}{enhanced,breakable,colback=sand,colframe=ink!20!white,boxrule=0.55pt,arc=1.4mm,left=3mm,right=3mm,top=2mm,bottom=2mm,before skip=7pt,after skip=8pt}
\newcommand{\anchorref}[1]{\par\smallskip{\footnotesize\textbf{Standard locator:} \texttt{#1}}\par}
\newcommand{\standard}[1]{\begin{standardbox}\textbf{Standard-grounded interpretation.} #1\end{standardbox}}
\newcommand{\practice}[1]{\begin{practicebox}\textbf{Engineering recommendation (not a normative ISO rule).} #1\end{practicebox}}
\newcommand{\takeaway}[1]{\par\textbf{Review question.} #1\par}
\newcommand{\docfront}{%
\begin{titlepage}
\vspace*{25mm}
{\color{teal}\bfseries\large ENGINEERING REFERENCE SERIES\par}
\vspace{12mm}
{\raggedright\color{ink}\bfseries\fontsize{24}{30}\selectfont\DocTitle\par}
\vspace{7mm}
{\color{muted}\fontsize{15}{19}\selectfont\DocSubtitle\par}
\vspace{12mm}
\noindent\color{teal}\rule{\linewidth}{2pt}\par
\vspace{9mm}
{\large\DocSummary\par}
\vfill
\begin{tabularx}{\linewidth}{@{}lX@{}}
\textbf{Volume} & \DocNumber\\
\textbf{Standards basis} & ISO/IEC 14882:2024 (seventh edition, C++23)\\
\textbf{Document type} & Original interpretive and instructional material\\
\textbf{Date} & 16 October 2026\\
\textbf{Toolchain} & pdfLaTeX; illustrative C++23 examples require separate compiler support\\
\end{tabularx}
\vspace{8mm}
{\footnotesize This document is not an authorized copy of, replacement for, or amendment to ISO/IEC 14882:2024. The original standard controls in the event of disagreement.}\par
\end{titlepage}
\pagenumbering{roman}
\tableofcontents
\clearpage
\pagenumbering{arabic}
\section*{How to use this volume}\addcontentsline{toc}{section}{How to use this volume}
The reference is \emph{ISO/IEC 14882:2024, Programming languages --- C++}. Section locators and square-bracketed stable names identify where to verify precise wording; they are not quotations. \textbf{Standard-grounded interpretation} boxes summarize language or library semantics. \textbf{Engineering recommendation} boxes contain optional design, security, portability, or process guidance. Code is original and illustrative; individual library facilities may not yet be implemented by every vendor. This suite does not grant rights to reproduce or distribute the licensed standard.\par\medskip
}
\newcommand{\docend}{\clearpage\section*{Reference and traceability}\addcontentsline{toc}{section}{Reference and traceability}
\textbf{Primary source.} ISO/IEC 14882:2024(en), \emph{Programming languages --- C++}, seventh edition, October 2024. Use its numbered clauses, stable names, Annex C (informative compatibility), and Annex D (compatibility features) for authoritative details.\par
\textbf{Limits.} Clause-level summaries and exercises in this document are selective explanations, not exhaustive reproductions of all semantic constraints, library mandates, complexity clauses, and overload participation requirements. A successful compiler test demonstrates behavior of that specific implementation and configuration, not universal conformance.\par
\textbf{Use of contemporary engineering tools.} Compilers, sanitizers, analyzers, CI systems, repository policies, and vulnerability taxonomies are implementation guidance rather than language-standard requirements.\par}

\begin{document}
\docfront
\section{Program shape and expected effort}
The plan is 28 instructional weeks at approximately 6--10 focused hours per week. Each module has explicit prerequisites; students may accelerate when they pass diagnostic entry assessments. The standards document is the source of semantics; toolchains, secure-development workflows and classroom scoring are additional educational design choices.

\small
\begin{longtable}{@{}p{.18\linewidth}p{.40\linewidth}p{.34\linewidth}@{}}
\toprule
\textbf{Track} & \textbf{Modules} & \textbf{Outcome}\\
\midrule
\endfirsthead
\toprule
\textbf{Track} & \textbf{Modules} & \textbf{Outcome}\\
\midrule
\endhead
Foundation & M01--M04 & Semantics, types, lifetime and initialization\\
\addlinespace[2pt]
Abstraction & M05--M08 & Exception-safe generic APIs and containers\\
\addlinespace[2pt]
Composition & M09--M12 & Ranges, modules, IO and concurrency\\
\addlinespace[2pt]
Assurance & M13--M16 & Portability, security, CI and capstone defense\\
\addlinespace[2pt]
\bottomrule
\end{longtable}
\normalsize
\section{Prerequisite graph}
The ordering is a directed acyclic dependency relation, not simply a numbering convention. An arrow means the earlier skill must be demonstrated before the later lab starts.

\begin{lstlisting}[language=]
M01 -> M02 -> M03 -> M05 -> M12
           |      |       |      |
           v      v       v      v
          M04 -> M06 -> M07 -> M08 -> M09
                    |       |       |
                    v       v       v
                   M10     M11 -> M13
                     \       |       /
                      \-> M14 -----/
                              |
                              v
                             M15 -> M16
\end{lstlisting}
\section{Module specifications}
\subsection{M01: Standards literacy and toolchain baselines}
\anchorref{1--4; 16}
Prerequisites: none. Duration: 2 week(s).

\textbf{Learning outcome.} Distinguish semantic classifications and prove selected toolchain language/library capabilities.

\textbf{Practical deliverable.} Write a diagnostic taxonomy; compile two small TUs; create environment manifest.

\practice{Exit gate: pass the knowledge check (80 percent), demonstrate the working lab, and cite the underlying clause(s) correctly.}
\subsection{M02: Lexical model, types and conversion rules}
\anchorref{5--7}
Prerequisites: M01. Duration: 2 week(s).

\textbf{Learning outcome.} Explain preprocessing, type identity, promotions, narrowing and value categories.

\textbf{Practical deliverable.} Build conversions matrix and compile-positive/negative examples.

\practice{Exit gate: pass the knowledge check (80 percent), demonstrate the working lab, and cite the underlying clause(s) correctly.}
\subsection{M03: Object lifetime and ownership}
\anchorref{6.7; 20; 23}
Prerequisites: M02. Duration: 2 week(s).

\textbf{Learning outcome.} Prove lifetime of owners, borrows and views across function boundaries.

\textbf{Practical deliverable.} Repair dangling view and invalid iterator exercises under ASan.

\practice{Exit gate: pass the knowledge check (80 percent), demonstrate the working lab, and cite the underlying clause(s) correctly.}
\subsection{M04: Control flow, declarations and initialization}
\anchorref{8--9}
Prerequisites: M02. Duration: 1 week(s).

\textbf{Learning outcome.} Choose initialization forms, evaluate loops, and trace constant evaluation.

\textbf{Practical deliverable.} Implement validated domain value types and static\_assert checks.

\practice{Exit gate: pass the knowledge check (80 percent), demonstrate the working lab, and cite the underlying clause(s) correctly.}
\subsection{M05: Class design and exception guarantees}
\anchorref{11; 14; 20}
Prerequisites: M03; M04. Duration: 2 week(s).

\textbf{Learning outcome.} Model class invariant, rule of zero, move behavior and exception recovery.

\textbf{Practical deliverable.} Implement guarded transaction with failure injection.

\practice{Exit gate: pass the knowledge check (80 percent), demonstrate the working lab, and cite the underlying clause(s) correctly.}
\subsection{M06: Overload resolution and callable interfaces}
\anchorref{7; 9; 12}
Prerequisites: M04. Duration: 1 week(s).

\textbf{Learning outcome.} Predict overload selection and identify ambiguous/different-conversion calls.

\textbf{Practical deliverable.} Build compile-time overload cases and annotated resolution reports.

\practice{Exit gate: pass the knowledge check (80 percent), demonstrate the working lab, and cite the underlying clause(s) correctly.}
\subsection{M07: Generic programming and concepts}
\anchorref{13; 18; 21}
Prerequisites: M05; M06. Duration: 2 week(s).

\textbf{Learning outcome.} Define a concept and test both syntax and semantic expectations.

\textbf{Practical deliverable.} Implement a generic accumulator with negative compile fixtures.

\practice{Exit gate: pass the knowledge check (80 percent), demonstrate the working lab, and cite the underlying clause(s) correctly.}
\subsection{M08: Strings, containers and iterators}
\anchorref{23--25}
Prerequisites: M03; M07. Duration: 2 week(s).

\textbf{Learning outcome.} Select storage models, reason about invalidation and choose iterator concepts.

\textbf{Practical deliverable.} Build a record index with iteration and erase invariants.

\practice{Exit gate: pass the knowledge check (80 percent), demonstrate the working lab, and cite the underlying clause(s) correctly.}
\subsection{M09: Ranges, algorithms and numerics}
\anchorref{26--28}
Prerequisites: M08. Duration: 2 week(s).

\textbf{Learning outcome.} Analyze lazy pipeline ownership, algorithm preconditions and performance.

\textbf{Practical deliverable.} Implement a range pipeline; compare eager and lazy evaluation.

\practice{Exit gate: pass the knowledge check (80 percent), demonstrate the working lab, and cite the underlying clause(s) correctly.}
\subsection{M10: Modules and migration engineering}
\anchorref{10; 15--16; Annex C}
Prerequisites: M07; M08. Duration: 1 week(s).

\textbf{Learning outcome.} Model module exports/reachability and plan mixed codebase migration.

\textbf{Practical deliverable.} Modularize an API when supported; publish fallback and build graph.

\practice{Exit gate: pass the knowledge check (80 percent), demonstrate the working lab, and cite the underlying clause(s) correctly.}
\subsection{M11: Error utilities, formatting and IO}
\anchorref{19; 22; 31--32}
Prerequisites: M05; M08. Duration: 2 week(s).

\textbf{Learning outcome.} Design structured errors, streams, path handling and input limits.

\textbf{Practical deliverable.} Implement strict parser with expected, size budget and unit tests.

\practice{Exit gate: pass the knowledge check (80 percent), demonstrate the working lab, and cite the underlying clause(s) correctly.}
\subsection{M12: Concurrency, atomics and coordination}
\anchorref{6.9; 33}
Prerequisites: M05; M09. Duration: 2 week(s).

\textbf{Learning outcome.} Prove shared-data synchronization and cancellation ownership.

\textbf{Practical deliverable.} Create a bounded concurrent queue; run race/liveness tests.

\practice{Exit gate: pass the knowledge check (80 percent), demonstrate the working lab, and cite the underlying clause(s) correctly.}
\subsection{M13: Security, static analysis and fuzzing}
\anchorref{6--33}
Prerequisites: M03; M09; M11; M12. Duration: 2 week(s).

\textbf{Learning outcome.} Classify vulnerabilities, design adversarial cases and interpret tool limitations.

\textbf{Practical deliverable.} Build fuzz and sanitizer harness; triage findings with provenance.

\practice{Exit gate: pass the knowledge check (80 percent), demonstrate the working lab, and cite the underlying clause(s) correctly.}
\subsection{M14: C++23 feature portability}
\anchorref{Annex C; 7--9; 22--31}
Prerequisites: M10; M11. Duration: 1 week(s).

\textbf{Learning outcome.} Adopt modern features based on evidence instead of flags alone.

\textbf{Practical deliverable.} Produce feature gates for expected, print, mdspan and library modules.

\practice{Exit gate: pass the knowledge check (80 percent), demonstrate the working lab, and cite the underlying clause(s) correctly.}
\subsection{M15: Integration, performance and release gates}
\anchorref{4; 16; 24--33}
Prerequisites: M12; M13; M14. Duration: 2 week(s).

\textbf{Learning outcome.} Combine functional, negative, instrumented, ABI and performance evidence.

\textbf{Practical deliverable.} Build a two-toolchain CI matrix and release evidence bundle.

\practice{Exit gate: pass the knowledge check (80 percent), demonstrate the working lab, and cite the underlying clause(s) correctly.}
\subsection{M16: Capstone assessment and standards defense}
\anchorref{1--33; Annex C--D}
Prerequisites: M15. Duration: 2 week(s).

\textbf{Learning outcome.} Defend design decisions with actual standard locators and measured outcomes.

\textbf{Practical deliverable.} Deliver secure bounded job processor, docs, tests, benchmark and review report.

\practice{Exit gate: pass the knowledge check (80 percent), demonstrate the working lab, and cite the underlying clause(s) correctly.}
\section{Assessment model}
\small
\begin{longtable}{@{}p{.18\linewidth}p{.40\linewidth}p{.34\linewidth}@{}}
\toprule
\textbf{Assessment} & \textbf{Weight} & \textbf{Acceptance rule}\\
\midrule
\endfirsthead
\toprule
\textbf{Assessment} & \textbf{Weight} & \textbf{Acceptance rule}\\
\midrule
\endhead
Conceptual quizzes & 20\% & At least 80\% across semantic vocabulary, clause interpretation and key distinctions\\
\addlinespace[2pt]
Engineering labs & 35\% & Tests pass; meaningful boundary cases; reviewable build\\
\addlinespace[2pt]
Standards defense & 20\% & Semantic claim mapped to cited clause and supported by reasoning\\
\addlinespace[2pt]
Capstone & 25\% & Architecture, security, correctness, portability and measurement evidence\\
\addlinespace[2pt]
\bottomrule
\end{longtable}
\normalsize
Passing the program requires at least 80 percent overall, at least 80 percent on the standards interpretation portion, and no unresolved critical memory-safety or data-race defect in the capstone. These are curriculum criteria, not ISO requirements.

\section{Capstone: secure bounded job processor}
Design a command-line library plus application that parses bounded job records and processes numeric workloads. Use validated value types, checked size arithmetic, expected-based recoverable errors, RAII, standard containers and algorithms. Add optional concurrency without violating ownership or cancellation contracts. Provide command-line, API and test documentation.

\small
\begin{longtable}{@{}p{.18\linewidth}p{.40\linewidth}p{.34\linewidth}@{}}
\toprule
\textbf{Milestone} & \textbf{Required artifact} & \textbf{Evidence}\\
\midrule
\endfirsthead
\toprule
\textbf{Milestone} & \textbf{Required artifact} & \textbf{Evidence}\\
\midrule
\endhead
Architecture & Module/namespace diagram and requirements & No hidden ownership or unbounded input path\\
\addlinespace[2pt]
Correctness & Core library + unit/property tests & Boundary and malformed inputs reproduced\\
\addlinespace[2pt]
Safety & Review checklist and sanitizer runs & Lifetime, bounds and concurrency reasoning\\
\addlinespace[2pt]
Portability & Feature probes and fallback strategy & Two documented toolchain configurations\\
\addlinespace[2pt]
Performance & Benchmark and profile & Recorded workload, p50/p95 and memory\\
\addlinespace[2pt]
Delivery & CI workflow, binary, report and changes & Build provenance, logs and reviewer sign-off\\
\addlinespace[2pt]
\bottomrule
\end{longtable}
\normalsize
\section{Prerequisite bypass examination}
\begin{itemize}[leftmargin=*,itemsep=1pt,topsep=3pt]
\item Explain why an integer overflow may have a different semantic status from a failed vector bounds precondition.
\item Compare a reference, unique\_ptr, shared\_ptr, span, and string\_view for ownership and absence semantics.
\item Predict overload choice from candidate viability and conversion ranks; identify ambiguity.
\item Explain why satisfied does not imply modeled for a standard concept.
\item Identify and fix a data race, then describe the synchronizes-with or mutex relationship.
\item Distinguish a compiler parsing C++23 syntax from a library shipping every new header.
\end{itemize}
\section{Instructor operating procedures}
\begin{itemize}[leftmargin=*,itemsep=1pt,topsep=3pt]
\item Begin each week with a concept diagnostic and a standards-locator challenge.
\item Record prerequisites completed and provide remediation before high-dependency modules.
\item Require each lab to ship source, toolchain identity, automated tests and a short design rationale.
\item Review unsafe fixtures only in isolated sandboxes and prohibit deployment of intentionally invalid examples.
\item Compare feature support against actual target environments before requiring new C++23 library APIs.
\item Close each module with a learner-authored summary and a brief oral standards defense.
\end{itemize}
\docend
\end{document}
